\documentclass[11pt]{article}
\usepackage{amsmath,amssymb,amsthm}
\newtheorem{theorem}{Theorem}
\begin{document}
\title{The next auxiliary-spin upper bound: third-order AC subspaces}
\author{}\date{}\maketitle
\begin{theorem}
If $N\ge6$ and $\mathcal C$ is a $k$-dimensional symmetric $N$-qubit
subspace of AC order at least three, then
\[
 \boxed{7k\le N+7.}
\]
This is an unrestricted upper bound. No reality, sparsity or symmetry
assumption on the subspace is made.
\end{theorem}
\begin{proof}
Put $j=N/2$, and let $S_a$ be the spin-three matrices on an auxiliary
seven-dimensional space. Set $X=\sum_aJ_a\otimes S_a$. If $P$ projects
onto the AC3 subspace, compression of a polynomial of degree at most
three in the physical spin matrices agrees with its normalized physical
trace times $P$.
This is a whole-code operator identity: polarization of the vanishing
rank-$1,2,3$ expectations gives $PT_{\ell m}P=0$ on every code
matrix element. Every ordered product of at most three $J_a$ decomposes
into tensor ranks at most three. Expanding $X^q$ as a sum of ordered
physical products tensored with ordered auxiliary products therefore
justifies the identities below on the full tensor-product code, including
arbitrary superpositions entangled with the auxiliary space. In particular,
\[
 (P\otimes I)X(P\otimes I)=0,
\]
\[
 (P\otimes I)X^2(P\otimes I)=4j(j+1)P\otimes I,
 \quad
 (P\otimes I)X^3(P\otimes I)=-2j(j+1)P\otimes I.
\]
The cubic identity follows from
$\operatorname{tr}(J_aJ_bJ_c)/(2j+1)=ij(j+1)\epsilon_{abc}/6$
and $\sum\epsilon_{abc}S_aS_bS_c=i\,3(3+1)I$.

Define
\begin{align*}
 F(X)&=(X+2j+5)(X+j+6)(X-2j+3),\\
 C_j&=15(2j^2+3j+6),\qquad W=F(X)-C_jI.
\end{align*}
Equivalently,
\[
 W=X^3+(j+14)X^2+(-4j^2+4j+63)X
      -2j(j+1)(2j+27)I.
\]
The three identities above show that $W$ has zero compression on
$\mathcal C\otimes V_3$.

In the total-spin sector $j+r$, $r=-3,-2,\ldots,3$, the eigenvalue of $X$
is
\[
 x_r=jr+\frac{r(r+1)}2-6.
\]
The corresponding eigenvalues of $W$ are
\[
\begin{array}{c|c}
 r& W(x_r)\\\hline
 -3&-5(j+1)(2j^2-3j+18)\\
 -2&-15(2j^2+3j+6)\\
 -1&-15(2j^2+3j+6)\\
 0&-2(2j+15)(j^2+j+3)\\
 1&-3(2j^3+15j^2+17j+30)\\
 2&-15(2j^2+3j+6)\\
 3&5j(4j^2+16j+27).
\end{array}
\]
All entries below the highest-spin sector are strictly negative, and the
highest entry is strictly positive. Indeed
$2j^2-3j+18=2(j-3/4)^2+135/8>0$.
Projection of the totally isotropic space $\mathcal C\otimes V_3$ into
the positive eigenspace is injective. The former has dimension $7k$,
and the latter is the spin-$(j+3)$ sector of dimension $N+7$.
\end{proof}

Combining this theorem with the rectangular-tensor bound gives
\[
 K(N,3)\le\min\left\{\left\lfloor\frac N6\right\rfloor,
                  \left\lfloor\frac{N+7}{7}\right\rfloor\right\},
 \qquad N\ge6,
\]
with the known exact saturation exclusions for the first term retained.
For example, at $N=55$ the old upper bound nine becomes eight.
No currently displayed catalogue entry (whose largest size is 42) is
improved by this third-order inequality alone.

This extends the already proved auxiliary-spin sequence
$3k\le N+3$ at order one and $5k\le N+5$ at order two.
It does not establish a general rule for higher orders, nor does it
establish attainment of the new upper bound.
\end{document}
